\ssrno{ЗАКЛЮЧЕНИЕ}

В результате выполнения выпускной квалификационной работы были разработаны метод обнаружения признаков генерации на изображении с использованием нейронной сети и его программная реализация.

Для достижения поставленной цели решены следующие задачи:
\begin{itemize}
	\item проанализирована предметная область;
	\item проанализированы известные нейросетевые решения;
	\item проведены классификация и сравнительный анализ рассмотренных решений;
	\item формализована задача в виде функциональной модели;
	\item спроектировано, разработано и протестировано программное обеспечение для обнаружения признаков генерации на изображении;
	\item проведено исследование работоспособности программной реализации предложенного метода.
\end{itemize}

В результате проведённого исследования установлено, что предложенный метод отвечает всем предъявленным требованиям: он превосходит по точности работы с немодифицированными изображениями все рассмотренные решения, кроме Capsule Forensics, обладая значением точности равным 98.75\% на наборе данных FaceForensics, и при этом превосходит Capsule Forensics по точности работы с низкокачественными изображениями, что является следствием его принадлежности к классу <<устойчивых>> методов.

Разработанный метод является подходящим в задачах, где критически важна максимальная точность при анализе изображений низкого качества, подверженных сжатию или иным видам зашумления. Применимость разработанного метода обусловлена повышенным требованием к точности и устойчивости при отсутствии ограничений вычислительной сложности. 

Дальнейшее развитие метода характеризуется повышением его точности на незашумленных изображениях набора данных FaceForensics, обобщением входных данных метода до произвольных изображений, а также уменьшением его вычислительной сложности при сохранении точности и устойчивости к шумам.
